\documentclass[a4paper,12pt]{article}
\usepackage[T2A]{fontenc}
\usepackage[utf8]{inputenc}
\usepackage[russian]{babel}
\usepackage{amsmath,amssymb,mathtools}
\usepackage{helvet}
\renewcommand{\familydefault}{\sfdefault}
\usepackage{icomma}
\usepackage[a4paper,left=22mm,right=22mm,top=23mm,bottom=23mm]{geometry}
\usepackage{microtype}
\usepackage{tikz}
\usetikzlibrary{arrows.meta,positioning,shapes.geometric,calc}
\usepackage{tkz-euclide}
\usepackage{booktabs,tabularx}
\usepackage{enumitem}
\usepackage{fancyhdr}
\usepackage{setspace}
\usepackage{xcolor}
\definecolor{accent}{HTML}{008080}
\definecolor{darkaccent}{HTML}{004D4D}
\definecolor{textgray}{HTML}{33383A}
\definecolor{softgray}{HTML}{667477}
\color{textgray}
\onehalfspacing
\setlength{\parindent}{0pt}
\setlength{\parskip}{4pt}
\setlist[itemize]{leftmargin=6mm,itemsep=3pt,topsep=3pt}
\setlist[enumerate]{leftmargin=8mm,itemsep=8pt,topsep=4pt}
\pagestyle{fancy}
\fancyhf{}
\fancyhead[L]{\footnotesize\color{softgray}Геометрия: свойства и вычисления}
\fancyhead[R]{\footnotesize\color{softgray}Кирилл Зиновьев}
\fancyfoot[C]{\footnotesize\color{softgray}\thepage}
\renewcommand{\headrulewidth}{0.3pt}
\setlength{\headheight}{16pt}
\newcommand{\mysection}[1]{\par\vspace{5pt}{\large\bfseries\color{darkaccent}#1}\par\vspace{1pt}{\color{accent}\hrule height .5pt}\vspace{7pt}}
\newcommand{\key}[1]{\textbf{\color{darkaccent}#1}}
\newcommand{\ang}[1]{\angle #1}
\begin{document}
\begin{titlepage}
\thispagestyle{empty}
\begin{center}
\vspace*{28mm}
{\large\color{darkaccent}МАТЕМАТИКА}\par\vspace{18mm}
{\LARGE Чек-лист}\par\vspace{12mm}
{\LARGE\color{accent}Медиана, биссектриса и высота:\\[2mm]
от свойства к уравнению}\par\vspace{18mm}
{\normalsize Геометрия. 7 класс}\par
\vfill
{\large Кирилл Зиновьев}\par\vspace{13mm}
{31 октября 2026 г.}\par\vspace{23mm}
{\small Лёвин Артём Александрович\\[2mm]
эксклюзивно для Кирилла Зиновьева}\par\vspace{18mm}
\begin{tikzpicture}[x=1cm,y=1cm]
\draw[accent,line width=.8pt] (0,0) .. controls (1.7,.70) and (3.1,-.58) .. (4.9,0)
 .. controls (6.8,.62) and (8.0,-.60) .. (9.8,0);
\end{tikzpicture}
\end{center}
\end{titlepage}
\setcounter{page}{1}

\mysection{1. Три свойства, на которых строится решение}
Сначала прочитайте условие и сделайте \key{чертёж}: подпишите вершины, точки на сторонах, известные длины и углы. Затем выпишите точное геометрическое свойство.

\key{Медиана} соединяет вершину с серединой противоположной стороны. Если $AM$ --- медиана $\triangle ABC$, то
\[
BM=MC=\frac{BC}{2},\qquad BC=BM+MC=2BM.
\]
\key{Биссектриса} делит угол на два равных угла. Если $BD$ --- биссектриса $\ang{ABC}$, то
\[
\ang{ABD}=\ang{DBC},\qquad
\ang{ABC}=\ang{ABD}+\ang{DBC}.
\]
\key{Высота} перпендикулярна прямой, содержащей противоположную сторону. Если $AD\perp BC$, то угол между $AD$ и $BC$ равен $90^\circ$.

\begin{center}
\begin{tikzpicture}[scale=.82]
\tkzDefPoint(0,0){B}
\tkzDefPoint(5.4,0){C}
\tkzDefPoint(1.7,3.2){A}
\tkzDefMidPoint(B,C)\tkzGetPoint{M}
\tkzDrawPolygon[line width=.9pt](A,B,C)
\tkzDrawSegment[accent,line width=1.3pt](A,M)
\tkzDrawPoints(A,B,C,M)
\tkzLabelPoint[above](A){$A$}
\tkzLabelPoints[below](B,M,C)
\tkzMarkSegments[mark=|,size=4pt](B,M M,C)
\end{tikzpicture}
\hspace{12mm}
\begin{tikzpicture}[scale=.82]
\tkzDefPoint(0,0){A}
\tkzDefPoint(5.8,0){C}
\tkzDefPoint(1.6,3.0){B}
\tkzDefLine[bisector](A,B,C)\tkzGetPoint{T}
\tkzInterLL(B,T)(A,C)\tkzGetPoint{D}
\tkzDrawPolygon[line width=.9pt](A,B,C)
\tkzDrawSegment[accent,line width=1.3pt](B,D)
\tkzDrawPoints(A,B,C,D)
\tkzLabelPoint[above](B){$B$}
\tkzLabelPoints[below](A,D,C)
\tkzMarkAngles[size=.65](A,B,D D,B,C)
\end{tikzpicture}
\end{center}
{\small\color{softgray}Слева отмечены равные части стороны; справа --- равные части угла. Эти признаки нельзя подменять друг другом.}

\mysection{2. Геометрическое свойство превращается в уравнение}
\key{Пример 1.} В $\triangle ABC$ проведена медиана $AM$; $BM=3x+2$, $MC=5x-8$. Найдите $x$ и $BC$.

По определению медианы $BM=MC$. Подставляем данные:
\[
3x+2=5x-8,\quad 2+8=5x-3x,\quad 10=2x,\quad x=5.
\]
Теперь находим \emph{обе} половины стороны:
\[
BM=3\cdot5+2=17,\qquad MC=5\cdot5-8=17.
\]
Значит, $BC=17+17=\boxed{34}$.

\key{Проверка по смыслу:} если $AM$ --- медиана, числа $BM$ и $MC$ обязаны совпасть. Если получились разные длины, ищите ошибку в подстановке или вычислении.

\clearpage
\mysection{3. Биссектриса и переменная в записи углов}
\key{Пример 2.} В $\triangle ABC$ отрезок $BD$ --- биссектриса $\ang{ABC}$, причём
\[
\ang{ABD}=4x-3,\qquad \ang{DBC}=2x+17.
\]
Найдите $x$ и $\ang{ABC}$.

Биссектриса делит угол пополам, поэтому
\[
4x-3=2x+17,\qquad 4x-2x=17+3,\qquad 2x=20,\quad x=10.
\]
Находим части угла и обязательно проверяем равенство:
\[
\ang{ABD}=4\cdot10-3=37^\circ,\quad
\ang{DBC}=2\cdot10+17=37^\circ.
\]
Весь угол равен сумме его частей:
\[
\boxed{\ang{ABC}=37^\circ+37^\circ=74^\circ}.
\]
\key{Важно:} равенство частей угла возникает из определения биссектрисы, а равенство частей стороны --- из определения медианы.

\mysection{4. Периметр: сначала формула}
\key{Пример 3.} В $\triangle ABC$ медиана $AM$, $AB=10$ см, $BC=12$ см, $AM=8$ см. Найдите периметр $\triangle ABM$.

Периметр --- сумма \emph{сторон нужного треугольника}:
\[
P_{ABM}=AB+BM+AM.
\]
Здесь $BM=BC:2=12:2=6$ см. Поэтому
\[
\boxed{P_{ABM}=10+6+8=24\text{ см}}.
\]
\key{Типичная ошибка:} вместо $BM$ подставить всю сторону $BC$. Следите за буквами в индексе периметра и проверяйте каждую сторону по чертежу.

\clearpage
\mysection{5. Когда один отрезок выполняет три роли}
В равнобедренном треугольнике медиана, проведённая к основанию, является одновременно биссектрисой и высотой. Верно и обратное: если медиана из вершины перпендикулярна противоположной стороне, то треугольник равнобедренный.

Если в $\triangle ABC$ точка $D$ лежит на $BC$ и
\[
BD=DC,\qquad \ang{BAD}=\ang{DAC},\qquad AD\perp BC,
\]
то $AD$ одновременно \key{медиана, биссектриса и высота}. Если $\ang{BAD}=32^\circ$, то $\ang{BAC}=64^\circ$.

\mysection{6. Самопроверка и границы вывода}
\begin{itemize}
\item После нахождения $x$ подставьте его \emph{во все} выражения из условия.
\item Проверьте, положительны ли найденные длины и согласуются ли равные части стороны или угла.
\item Если нужна формульная величина, сначала запишите её формулу: так станет ясно, каких данных не хватает.
\item Не приписывайте медиане свойства биссектрисы или высоты, если они не даны и не доказаны.
\item Если данных недостаточно, укажите это и назовите недостающую величину.
\end{itemize}

\clearpage
\thispagestyle{fancy}
\begin{center}
{\Large\bfseries\color{darkaccent}Алгоритм: от чертежа к ответу}\par\vspace{9mm}
{\small\color{softgray}Единый маршрут для вычислительных задач с медианой, биссектрисой и высотой}
\end{center}
\vspace{9mm}
\begin{center}
\begin{tikzpicture}[
  >=Stealth,
  s/.style={ellipse,draw=darkaccent,line width=.8pt,align=center,text width=5.6cm,minimum height=1cm,font=\small},
  a/.style={rounded corners=2pt,rectangle,draw=accent,line width=.8pt,align=center,text width=8.3cm,minimum height=1.15cm,inner sep=5pt,font=\small},
  d/.style={diamond,aspect=2.6,draw=darkaccent,line width=.8pt,align=center,text width=4cm,inner sep=1.6pt,font=\small},
  side/.style={rectangle,rounded corners=2pt,draw=accent,line width=.7pt,align=center,text width=3.4cm,minimum height=1.5cm,inner sep=4pt,font=\small},
  arr/.style={-Stealth,line width=.8pt,color=darkaccent},
  node distance=10mm]
\node[s] (start) {Начало задачи};
\node[a,below=of start] (drawing) {Сделайте чертёж, отметьте известные данные и нужный отрезок или угол};
\node[a,below=of drawing] (theory) {Запишите свойство медианы, биссектрисы или высоты; если нужно, формулу периметра};
\node[d,below=13mm of theory] (q) {Известны все величины для вычисления?};
\node[side,right=21mm of q] (solve) {Нет: составьте уравнение и найдите $x$};
\node[a,below=20mm of q] (calc) {Подставьте найденные значения; вычислите нужную величину};
\node[a,below=of calc] (check) {Проверьте равенства, положительность длин и смысл результата};
\node[s,below=of check] (end) {Запишите ответ};
\foreach \f/\t in {start/drawing,drawing/theory,theory/q,calc/check,check/end}{\draw[arr] (\f.south)--(\t.north);}
\draw[arr] (q.south)--node[right,font=\footnotesize]{да}(calc.north);
\draw[arr] (q.east)--node[above,font=\footnotesize]{нет}(solve.west);
\draw[arr] (solve.south) |- ([yshift=8mm]calc.east) -- (calc.east);
\end{tikzpicture}
\end{center}

\clearpage
\mysection{Тренировочный блок --- Путь А}
\textbf{10 заданий.} Выполняйте их по порядку. В расчётных задачах сделайте чертёж, запишите свойство и только после этого составляйте уравнение или формулу.
\begin{enumerate}[label=\textbf{\color{accent}\arabic*.},itemsep=12pt]
\item В $\triangle ABC$ отрезок $AM$ --- медиана, $BC=22$ см. Найдите $BM$ и $MC$.
\item В $\triangle ABC$ отрезок $BD$ --- биссектриса $\ang{ABC}$, равного $86^\circ$. Найдите $\ang{ABD}$ и $\ang{DBC}$.
\item В $\triangle ABC$ проведена медиана $AM$, $BM=3x+4$, $MC=5x-8$. Найдите $x$ и длину $BC$.
\item В $\triangle ABC$ отрезок $BD$ --- биссектриса $\ang{ABC}$, $\ang{ABD}=4x-7$, $\ang{DBC}=2x+17$. Найдите $x$ и $\ang{ABC}$.
\item В $\triangle ABC$ медиана $AM$, $AB=12$ см, $BC=18$ см, $AM=10$ см. Найдите периметр $\triangle ABM$.
\item В $\triangle ABC$ проведена медиана $BK$, $AB=15$ см, $AC=20$ см, $BK=8$ см. Найдите периметр $\triangle ABK$.
\item В $\triangle ABC$ медиана $AM$, $BM=4x+1$, $MC=6x-7$. При $x=4$ записали: $BM=17$, $MC=5$, $BC=22$. Найдите ошибку и правильную длину $BC$.
\item В $\triangle ABC$ медиана $AM$, $AB=10$ см, $BC=12$ см. Можно ли однозначно вычислить периметр $\triangle ABM$? Объясните, какой величины не хватает.
\item В $\triangle ABC$ точка $D$ лежит на стороне $BC$, $BD=DC$, $\ang{BAD}=\ang{DAC}$ и $AD\perp BC$. Чем является $AD$? Найдите $\ang{BAC}$, если $\ang{BAD}=28^\circ$.
\item В $\triangle ABC$ медиана $AM$, $BC=2x+14$, $BM=3x-1$, $AB=12$ см, $AM=9$ см. Найдите $x$, $BC$ и периметр $\triangle ABM$.
\end{enumerate}

\clearpage
\mysection{Ответы к тренировочному блоку}
\renewcommand{\arraystretch}{1.45}
\begin{tabularx}{\linewidth}{@{}>{\bfseries}p{10mm}X@{}}
\toprule
№ & Ответ или контрольный вывод \\
\midrule
1 & $BM=MC=11$ см. \\
2 & $\ang{ABD}=\ang{DBC}=43^\circ$. \\
3 & $x=6$; $BC=44$. \\
4 & $x=12$; $\ang{ABC}=82^\circ$. \\
5 & $P_{ABM}=31$ см. \\
6 & $P_{ABK}=33$ см. \\
7 & Неверно вычислен $MC$: $6\cdot4-7=17$. Поэтому $BC=34$. \\
8 & Нет: известно $BM=6$ см, но длина $AM$ не задана и не определяется указанными данными. \\
9 & $AD$ --- медиана, биссектриса и высота; $\ang{BAC}=56^\circ$. \\
10 & $x=4$; $BC=22$ см; $P_{ABM}=32$ см. \\
\bottomrule
\end{tabularx}
\vspace{13mm}
\mysection{Памятка перед сдачей}
\begin{tabularx}{\linewidth}{@{}p{45mm}X@{}}
\toprule
Признак в условии & Что можно сразу записать \\
\midrule
Медиана $AM$ & $BM=MC=BC/2$. \\
Биссектриса $BD$ & $\ang{ABD}=\ang{DBC}$. \\
Высота $AD$ & $AD\perp BC$ (или прямой, содержащей $BC$). \\
Периметр $\triangle ABM$ & $P_{ABM}=AB+BM+AM$. \\
\bottomrule
\end{tabularx}
\end{document}
